\section{Conditional low-acceleration matter coupling}

\subsection{Inverse construction of the cubic-gradient operator}

Nonlinear Poisson equations derived from nonquadratic potential actions have a
standard precedent in AQUAL-type theories \cite{bekensteinmilgrom1984}.  Here
we isolate the coefficient normalization required by the ITSM infrared branch.
Let $q=|\nabla\psi|$ and consider the static force-scalar action
\begin{equation}
 S_\psi=\int \dd t\,\dd^3x\left[-K(q)-C_m\rho_b\psi\right].
 \label{eq:static-general}
\end{equation}
Variation yields
\begin{equation}
 \nabla\cdot\left[\frac{K'(q)}{q}\nabla\psi\right]=C_m\rho_b.
 \label{eq:general-field}
\end{equation}
The required cubic-gradient operator is parameterized as
\begin{equation}
 K(q)=\frac{C_{\rm IR}}{12\pi Ga_0}q^3,
 \qquad C_{\rm IR}>0.
 \label{eq:cubic-K}
\end{equation}

The complete static weak-field action in this normalization is
\begin{align}
 S_{\rm WF}=\int \dd t\,\dd^3x\bigg[
 &-\frac{|\nabla\Phi_N|^2}{8\pi G}-\rho_b\Phi_N \nonumber\\
 &-\frac{C_{\rm IR}}{12\pi Ga_0}|\nabla\psi|^3
 -C_m\rho_b\psi\bigg].
 \label{eq:wf-action}
\end{align}
For independent positive coefficients $C_m$ and $C_{\rm IR}$, the sourced
equation is
\begin{equation}
 \nabla\cdot\left[\frac{|\nabla\psi|}{a_0}\nabla\psi\right]
 =4\pi G\frac{C_m}{C_{\rm IR}}\rho_b.
 \label{eq:p3-local}
\end{equation}
For spherical symmetry this gives
\begin{equation}
 |\nabla\psi|^2=\frac{C_m}{C_{\rm IR}}a_0g_N.
 \label{eq:general-gradient}
\end{equation}
Since the physical acceleration is $g_P=C_m|\nabla\psi|$, the invariant
observable square-root coefficient is
\begin{equation}
 g_P=C_{\rm obs}\sqrt{a_0g_N},
 \qquad
 \boxed{C_{\rm obs}=\frac{C_m^{3/2}}{\sqrt{C_{\rm IR}}}}.
 \label{eq:observable-coupling}
\end{equation}
Only under the normalization choice $C_m=C_{\rm IR}=C$ does
$C_{\rm obs}=C$.  That convenient choice is not a physical derivation of
$C=2/3$; MAT-001 must calculate the invariant combination in
Eq.~\eqref{eq:observable-coupling}.

\subsection{Compact toroidal source}

On compact $\Tthree$ there is no boundary at spatial infinity, and a periodic
divergence integrates to zero.  Write
\begin{equation}
 \rho_b=\bar\rho_b+\delta\rho_b,
 \qquad
 \psi=\bar\psi(t)+\varphi(t,\mathbf x),
\end{equation}
with zero spatial averages for $\delta\rho_b$ and $\varphi$.  The local
periodic equation has the form
\begin{equation}
 D_i\left[\frac{|D\varphi|}{a_0}D^i\varphi\right]
 =4\pi G\frac{C_m}{C_{\rm IR}}\delta\rho_b+\mathcal S_Q,
 \qquad
 \int_{\Tthree}\left(4\pi G\frac{C_m}{C_{\rm IR}}
 \delta\rho_b+\mathcal S_Q\right)\,\dd V=0.
 \label{eq:periodic-p3}
\end{equation}
Here $\mathcal S_Q\equiv\mathcal S_Q[\delta Q_{\rm syn}]$ is a conditional
zero-mean scalar source representing the inhomogeneous projection of
reservoir--plenum exchange into the force-scalar equation.  Its functional map from
the vector exchange current has not been derived; setting $\mathcal S_Q=0$ is
the local-static baseline rather than a claim that reservoir perturbations
vanish identically.
The homogeneous mode belongs to the cosmological/open-system problem.  A
localized spherical solution is therefore the small-region,
compensated-overdensity limit of Eq.~\eqref{eq:periodic-p3}, not an isolated
global cosmology.

\subsection{Disk geometry and coefficient status}

For the flat periodic metric and local-static baseline $\mathcal S_Q=0$,
equality of divergences permits the torus Hodge decomposition
\begin{equation}
 \frac{|\nabla\varphi|}{a_0}\nabla\varphi
 =\frac{C_m}{C_{\rm IR}}\nabla\Phi_N+\nabla\times\mathbf H
 +\mathbf h_0.
 \label{eq:curl-field}
\end{equation}
Here $\mathbf H$ is periodic and $\mathbf h_0$ is the constant harmonic flux,
equal to the spatial mean of the difference of the two gradient-flux terms.
The mean of a periodic curl vanishes, whereas the mean of
$|\nabla\varphi|\nabla\varphi$ need not vanish even when
$\langle\nabla\varphi\rangle=0$.  Thus $\mathbf h_0=0$ requires an additional
condition; it does not follow from periodicity alone.  This flux is determined
by the solution, not a freely adjustable uniform acceleration or a winding
of the force potential.  A curved spatial metric requires its corresponding
Hodge decomposition.  For $\mathcal S_Q\ne0$, its separately derived source
contribution must also be subtracted before taking a divergence-free difference.

Sufficient local symmetry and boundary conditions can remove the solenoidal
correction, but global spherical symmetry is not a property of $\Tthree$.
For a finite disk the correction is generally nonzero.
Consequently a pointwise square-root prescription is a
historical algebraic approximation.  DISK-001 must solve
Eq.~\eqref{eq:periodic-p3} for realistic baryonic sources before
$C_{\rm obs}=2/3$ and $C_{\rm obs}\simeq1$ can be fairly compared.

The separate coefficients change under a rescaling of $\psi$, while the
physical response is carried by the invariant combination
$C_{\rm obs}=C_m^{3/2}/\sqrt{C_{\rm IR}}$.  The trace ratio $2/3$ remains a
geometric motivation and a matching hypothesis, not a completed local Wilson
coefficient.
